% Banque d'exercices — chapitre 9
% Le chapitre est déterminé par le nom de ce fichier.

\begin{exo}[Relation de Mayer : $C_p - C_v = nR$]{mayer-relation}
\keywords{coefficients calorimétriques, Cp, Cv, relation de Mayer, gaz parfait, coefficient de dilatation}

On considère un gaz parfait de $n$ moles. On rappelle que $C_p = \left(\partial H/\partial T\right)_P$ et $C_v = \left(\partial U/\partial T\right)_V$, et que $H = U + PV$.

\begin{enumerate}
  \item Exprimer $H$ en fonction de $U$ et montrer que $H = U + nRT$ pour un gaz parfait.
  \item En déduire que $C_p - C_v = nR$ (relation de Mayer).
  \item Application numérique : pour l'air ($n=1$ mol, $R = 8{,}314\,\text{J·mol}^{-1}\text{K}^{-1}$), calculer $C_p - C_v$.
  \item Montrer que la relation de Mayer implique $\gamma = C_p/C_v > 1$.
\end{enumerate}

\begin{indication}
Utiliser $PV = nRT$ pour simplifier $H = U + PV$, puis dériver par rapport à $T$ à pression constante.
\end{indication}

\begin{solution}
\textbf{1.} $H = U + PV = U + nRT$ (gaz parfait). Comme $U$ ne dépend que de $T$, $H$ ne dépend aussi que de $T$.

\textbf{2.} $C_p = \left(\partial H/\partial T\right)_P = \mathrm{d}H/\mathrm{d}T = \mathrm{d}U/\mathrm{d}T + nR = C_v + nR$. Donc $C_p - C_v = nR$.

\textbf{3.} $C_p - C_v = 1 \times 8{,}314 \approx 8{,}3\,\text{J·K}^{-1}$.

\textbf{4.} $C_p = C_v + nR > C_v$ car $nR > 0$. Donc $\gamma = C_p/C_v > 1$.
\end{solution}

\end{exo}

\begin{exo}[Les six coefficients du gaz parfait]{coefficients-gaz-parfait}
\keywords{coefficients calorimétriques, coefficients thermoélastiques, dilatation isobare, compressibilité isotherme, gaz parfait}

On rappelle les définitions, pour $n$ moles de gaz parfait :
\[
\delta Q = C_V\,\mathrm{d}T + \ell\,\mathrm{d}V = C_P\,\mathrm{d}T + h\,\mathrm{d}P, \qquad
\alpha = \frac{1}{V}\left(\frac{\partial V}{\partial T}\right)_P, \quad
\beta = \frac{1}{P}\left(\frac{\partial P}{\partial T}\right)_V, \quad
\chi_T = -\frac{1}{V}\left(\frac{\partial V}{\partial P}\right)_T.
\]

\begin{enumerate}
\item En utilisant le fait que $U$ ne dépend que de $T$, montrer que $\ell = P$ pour le gaz parfait. En utilisant le fait que $H$ ne dépend que de $T$, montrer que $h = -V$.
\item Calculer $\alpha$, $\beta$ et $\chi_T$ pour le gaz parfait.
\item Vérifier sur ces valeurs les relations générales $\ell = \beta TP$, $h = -\alpha VT$ et $P\beta\chi_T = \alpha$.
\item Démontrer la relation $\ell = h\left(\dfrac{\partial P}{\partial V}\right)_T$ et la vérifier pour le gaz parfait.
\item Interpréter physiquement le signe de $h$ : que doit faire la pression d'un gaz parfait qui reçoit de la chaleur à température constante ?
\end{enumerate}

\begin{indication}
Isotherme : $\mathrm{d}U = 0$ donne $\delta Q = P\,\mathrm{d}V$ ; de même $\mathrm{d}H = 0$ donne $\delta Q = -V\,\mathrm{d}P$. Pour la question 4, écrire $\delta Q$ en variables $(T,V)$ et y substituer $\mathrm{d}P$ exprimé en fonction de $\mathrm{d}T$ et $\mathrm{d}V$... ou plus simplement comparer les deux écritures de $\delta Q$ sur une isotherme.
\end{indication}

\begin{solution}
\textbf{Question 1.} Sur une isotherme du gaz parfait, $\mathrm{d}U = 0$, donc $\delta Q = P\,\mathrm{d}V$ ; par identification avec $\delta Q = C_V\,\mathrm{d}T + \ell\,\mathrm{d}V = \ell\,\mathrm{d}V$, on lit $\ell = P$. De même $H$ ne dépend que de $T$ : $\mathrm{d}H = \delta Q + V\,\mathrm{d}P = 0$ sur l'isotherme, donc $\delta Q = -V\,\mathrm{d}P = h\,\mathrm{d}P$, d'où $h = -V$.

\textbf{Question 2.} De $V = nRT/P$ : $\alpha = \dfrac{1}{V}\cdot\dfrac{nR}{P} = \dfrac{1}{T}$ et $\chi_T = -\dfrac{1}{V}\left(-\dfrac{nRT}{P^2}\right) = \dfrac{1}{P}$. De $P = nRT/V$ : $\beta = \dfrac{1}{P}\cdot\dfrac{nR}{V} = \dfrac{1}{T}$.

\textbf{Question 3.} $\beta TP = \frac{1}{T}\,TP = P = \ell$ \checkmark ; $-\alpha VT = -\frac{1}{T}\,VT = -V = h$ \checkmark ; $P\beta\chi_T = P\cdot\frac{1}{T}\cdot\frac{1}{P} = \frac{1}{T} = \alpha$ \checkmark.

\textbf{Question 4.} Sur une isotherme, les deux écritures de $\delta Q$ donnent $\ell\,\mathrm{d}V = h\,\mathrm{d}P$, et comme $\mathrm{d}P = \left(\partial P/\partial V\right)_T\mathrm{d}V$ à $T$ fixé :
\[
\ell = h\left(\frac{\partial P}{\partial V}\right)_T.
\]
Gaz parfait : $h\left(\partial P/\partial V\right)_T = (-V)\left(-\dfrac{nRT}{V^2}\right) = \dfrac{nRT}{V} = P = \ell$. \checkmark

\textbf{Question 5.} $\delta Q = h\,\mathrm{d}P = -V\,\mathrm{d}P$ : pour recevoir de la chaleur ($\delta Q > 0$) à température constante, il faut $\mathrm{d}P < 0$. Rien d'étrange : à $T$ fixée, $PV$ est constant, le gaz se dilate en absorbant la chaleur ($\delta Q = P\,\mathrm{d}V > 0$) et sa pression chute d'autant.
\end{solution}
\end{exo}

\begin{exo}[Entropie du CO$_2$ : quand $C_V$ dépend de $T$]{cv-dependant-de-t}
\keywords{capacité thermique variable, dioxyde de carbone, entropie, chauffage isochore, intégration}

Entre 273 et 500 K, la capacité thermique molaire à volume constant du dioxyde de carbone est bien décrite par
\[
c_V(T) = A + BT, \qquad A = 23{,}83\ \text{J·mol}^{-1}\text{K}^{-1}, \quad B = 22{,}15\times10^{-3}\ \text{J·mol}^{-1}\text{K}^{-2}.
\]
On traite le CO$_2$ comme un gaz parfait (au sens $PV = nRT$).

\begin{enumerate}
\item Établir, à une constante près, l'expression de l'entropie de $n$ moles de ce gaz en fonction de $T$ et $V$.
\item Calculer la variation d'entropie $\Delta S$ pour un échauffement isochore de cinq moles de CO$_2$ de 289 K à 400 K.
\item La relation de Mayer molaire $c_P - c_V = R$ reste-t-elle valable ? Que vaut $c_P(T)$ ?
\end{enumerate}

\begin{indication}
Reprendre $\mathrm{d}S = \dfrac{nc_V(T)}{T}\mathrm{d}T + \dfrac{nR}{V}\mathrm{d}V$ et intégrer terme à terme ; le terme en $BT/T$ donne une contribution linéaire en $T$.
\end{indication}

\begin{solution}
\textbf{Question 1.}
\[
S(T,V) = n\!\int\!\frac{A + BT}{T}\,\mathrm{d}T + nR\ln V + \text{cste} = n\left[A\ln T + BT\right] + nR\ln V + \text{cste}.
\]

\textbf{Question 2.} Isochore, $n = 5$ mol :
\[
\Delta S = 5\left[23{,}83\ln\frac{400}{289} + 22{,}15\times10^{-3}\,(400 - 289)\right] = 5\left[7{,}75 + 2{,}46\right] \approx 51\ \text{J/K}.
\]
Le terme linéaire en $T$ n'est pas négligeable : il pèse près d'un quart du total.

\textbf{Question 3.} Oui : la relation de Mayer du gaz parfait ne dépend que de l'équation d'état $PV = nRT$, pas de la forme de $c_V$. Donc $c_P(T) = c_V(T) + R = (A + R) + BT$ : les deux capacités dépendent de $T$, mais leur différence reste constante.
\end{solution}
\end{exo}

\begin{exo}[Relations de Clapeyron : la loi de Joule démontrée]{relations-clapeyron}
\keywords{relations de Clapeyron, coefficients calorimétriques, relation de Maxwell, loi de Joule, énergie interne, gaz parfait}

On rappelle les relations de Maxwell issues de $F$ et de $G$ :
\[
\left(\frac{\partial S}{\partial V}\right)_T = \left(\frac{\partial P}{\partial T}\right)_V, \qquad
\left(\frac{\partial S}{\partial P}\right)_T = -\left(\frac{\partial V}{\partial T}\right)_P.
\]

\begin{enumerate}
\item À partir des définitions $\ell = T\left(\partial S/\partial V\right)_T$ et $h = T\left(\partial S/\partial P\right)_T$, démontrer les deux relations de Clapeyron :
\[
\ell = T\left(\frac{\partial P}{\partial T}\right)_V, \qquad h = -T\left(\frac{\partial V}{\partial T}\right)_P.
\]
\item En déduire l'identité générale
\[
\left(\frac{\partial U}{\partial V}\right)_T = T\left(\frac{\partial P}{\partial T}\right)_V - P.
\]
\item Montrer que pour tout fluide dont la pression est proportionnelle à $T$ à volume fixé, l'énergie interne ne dépend pas du volume. Vérifier explicitement pour le gaz parfait : c'est la (première) loi de Joule, ici \emph{démontrée} et non postulée.
\item L'expérience de la détente de Joule montre en réalité un léger refroidissement pour les gaz réels. Qu'est-ce que cela révèle sur $\left(\partial P/\partial T\right)_V$ des gaz réels ?
\end{enumerate}

\begin{indication}
Pour la question 2, partir de $\mathrm{d}U = T\,\mathrm{d}S - P\,\mathrm{d}V$ et dériver par rapport à $V$ à $T$ fixé. Pour la question 3, écrire $P = T\,f(V)$ et calculer.
\end{indication}

\begin{solution}
\textbf{Question 1.} Immédiat en substituant les relations de Maxwell dans les définitions : $\ell = T\left(\partial S/\partial V\right)_T = T\left(\partial P/\partial T\right)_V$, et $h = T\left(\partial S/\partial P\right)_T = -T\left(\partial V/\partial T\right)_P$.

\textbf{Question 2.} De $\mathrm{d}U = T\,\mathrm{d}S - P\,\mathrm{d}V$, à $T$ fixé :
\[
\left(\frac{\partial U}{\partial V}\right)_T = T\left(\frac{\partial S}{\partial V}\right)_T - P = \ell - P = T\left(\frac{\partial P}{\partial T}\right)_V - P.
\]

\textbf{Question 3.} Si $P = T\,f(V)$, alors $T\left(\partial P/\partial T\right)_V = T\,f(V) = P$, et $\left(\partial U/\partial V\right)_T = 0$ : $U$ ne dépend que de $T$. Le gaz parfait ($P = nRT/V$) est exactement de cette forme : $T\cdot nR/V - P = 0$. \checkmark\ La loi de Joule n'est donc pas une hypothèse supplémentaire : elle découle de l'équation d'état et des deux principes.

\textbf{Question 4.} Un refroidissement dans la détente de Joule signifie $\left(\partial U/\partial V\right)_T > 0$ ($U$ augmente avec $V$ à $T$ fixé, donc à $U$ constant, $T$ baisse quand $V$ croît) : donc $T\left(\partial P/\partial T\right)_V > P$ pour un gaz réel. La pression y croît plus vite que linéairement en température — signature des interactions attractives entre molécules, comme le montre l'exercice suivant sur van der Waals.
\end{solution}
\end{exo}

\begin{exo}[Gaz de van der Waals : énergie interne et détente de Joule]{van-der-waals-energie}
\keywords{van der Waals, gaz réel, énergie interne, détente de Joule, refroidissement, pression de cohésion}

Le gaz de van der Waals a pour équation d'état
\[
P = \frac{nRT}{V - nb} - \frac{an^2}{V^2},
\]
avec $a, b > 0$. On suppose $C_V$ constant.

\begin{enumerate}
\item Calculer $\ell = T\left(\partial P/\partial T\right)_V$ puis $\left(\partial U/\partial V\right)_T$.
\item En déduire que $U(T,V) = C_VT - \dfrac{an^2}{V} + \text{cste}$. Interpréter le signe du terme en $a$.
\item Ce gaz subit une détente de Joule (récipient calorifugé, détente dans le vide) de $V_0$ vers $2V_0$. Montrer que
\[
\Delta T = -\frac{an^2}{2V_0\,C_V} < 0.
\]
\item Application numérique pour une mole de diazote : $a = 0{,}137$ Pa·m$^6$·mol$^{-2}$, $C_V = \frac{5}{2}R$, $V_0 = 1{,}0$ L. Commenter la petitesse du résultat et expliquer pourquoi l'expérience historique de Joule (avec un calorimètre à eau) n'a rien détecté.
\end{enumerate}

\begin{indication}
Utiliser l'identité $\left(\partial U/\partial V\right)_T = T\left(\partial P/\partial T\right)_V - P$ de l'exercice précédent, puis intégrer à $T$ fixé. Dans la détente de Joule, $\Delta U = 0$.
\end{indication}

\begin{solution}
\textbf{Question 1.} $\left(\dfrac{\partial P}{\partial T}\right)_V = \dfrac{nR}{V - nb}$, donc $\ell = \dfrac{nRT}{V - nb} = P + \dfrac{an^2}{V^2}$ et
\[
\left(\frac{\partial U}{\partial V}\right)_T = \ell - P = \frac{an^2}{V^2} > 0.
\]

\textbf{Question 2.} En intégrant à $T$ fixé : $U = -\dfrac{an^2}{V} + f(T)$, et $C_V = f'(T)$ constant donne $f = C_VT + \text{cste}$ :
\[
U(T,V) = C_VT - \frac{an^2}{V} + \text{cste}.
\]
Le terme $-an^2/V$ est une énergie potentielle d'attraction entre molécules : rapprocher les molécules (diminuer $V$) abaisse cette énergie, et il faut fournir du travail contre l'attraction pour dilater le gaz.

\textbf{Question 3.} Détente de Joule : $W = Q = 0$ donc $\Delta U = 0$ :
\[
C_V\Delta T - an^2\left(\frac{1}{2V_0} - \frac{1}{V_0}\right) = 0 \implies \Delta T = -\frac{an^2}{2V_0C_V} < 0.
\]
L'énergie d'agitation thermique se convertit en énergie potentielle d'attraction : le gaz se refroidit — contrairement au gaz parfait dont la température ne bouge pas.

\textbf{Question 4.} $\Delta T = -\dfrac{0{,}137}{2\times1{,}0\times10^{-3}\times20{,}8} \approx -3{,}3$ K. C'est faible, et dans l'expérience historique la capacité thermique du calorimètre à eau (des kilojoules par kelvin, contre $\sim21$ J/K pour le gaz) diluait cet effet d'un facteur cent : Joule ne pouvait pas le voir. Les mesures modernes, directes sur le gaz, le confirment.
\end{solution}
\end{exo}

\begin{exo}[Dilatation thermique des océans]{dilatation-oceans}
\keywords{coefficient de dilatation isobare, océans, réchauffement climatique, montée du niveau de la mer, capacité thermique de l'eau}

Les rapports du GIEC estiment que l'excès de flux radiatif dû à l'effet de serre injecte dans les océans environ $E = 5\times10^{21}$ J par an. Données : coefficient de dilatation isobare de l'eau $\alpha \approx 2{,}5\times10^{-4}\ \text{K}^{-1}$ (à 20 °C), capacité thermique massique $c \approx 4180$ J·kg$^{-1}$·K$^{-1}$, masse volumique $\rho = 1000$ kg·m$^{-3}$, profondeur moyenne des océans $H = 4000$ m, rayon terrestre $R_T = 6400$ km, fraction de la surface terrestre couverte par les océans : $75\,\%$.

\begin{enumerate}
\item Convertir $E$ en puissance moyenne. Comparer à la puissance d'une centrale nucléaire ($\sim10^9$ W).
\item Estimer la masse totale des océans, puis l'élévation annuelle de température $\Delta T$ qui résulte de l'apport $E$.
\item Montrer qu'une colonne d'eau de hauteur $H$ chauffée de $\Delta T$ s'élève de $\Delta h = \alpha H\,\Delta T$. Estimer la montée annuelle du niveau de la mer par simple dilatation thermique, puis sur un siècle.
\item Citer au moins une raison pour laquelle la hausse réellement observée ($\sim3$ mm/an) est plus grande.
\end{enumerate}

\begin{indication}
Une année vaut environ $3{,}15\times10^7$ s. Pour la question 3, le volume de la colonne de section $A$ vaut $V = AH$ et sa variation isobare est $\Delta V = \alpha V\Delta T$ ; la section ne changeant pas, tout passe dans la hauteur.
\end{indication}

\begin{solution}
\textbf{Question 1.} $\mathcal{P} = \dfrac{5\times10^{21}}{3{,}15\times10^7} \approx 1{,}6\times10^{14}$ W, soit l'équivalent d'environ 160\,000 centrales nucléaires fonctionnant en continu.

\textbf{Question 2.} Surface océanique : $S \approx 0{,}75\times4\pi R_T^2 \approx 3{,}9\times10^{14}$ m$^2$, d'où une masse $m = \rho SH \approx 1{,}5\times10^{21}$ kg. Alors
\[
\Delta T = \frac{E}{mc} \approx \frac{5\times10^{21}}{1{,}5\times10^{21}\times4180} \approx 8\times10^{-4}\ \text{K/an},
\]
soit environ un millième de degré par an (0,1 °C par siècle en ordre de grandeur).

\textbf{Question 3.} $\Delta V = \alpha V\Delta T$ avec $V = AH$ et $A$ fixée : $\Delta h = \alpha H\Delta T \approx 2{,}5\times10^{-4}\times4000\times8\times10^{-4} \approx 0{,}8$ mm/an, soit de l'ordre de 8 cm par siècle par pure dilatation thermique.

\textbf{Question 4.} La dilatation thermique n'est qu'une contribution ($\sim$ un tiers) : il faut ajouter la fonte des glaces continentales (Groenland, Antarctique, glaciers), et tenir compte du fait que $\alpha$ croît avec la température (le chauffage se concentre dans les couches superficielles plus chaudes, qui se dilatent davantage que la moyenne).
\end{solution}
\end{exo}

\begin{exo}[Relation de Mayer générale et stabilité : $\gamma > 1$]{mayer-generale-stabilite}
\keywords{relation de Mayer générale, stabilité thermodynamique, gamma, compressibilité, dilatation, capacités thermiques}

On rappelle la relation de Mayer générale établie en cours :
\[
C_P - C_V = T\left(\frac{\partial P}{\partial T}\right)_V\left(\frac{\partial V}{\partial T}\right)_P,
\]
ainsi que les définitions de $\alpha$, $\beta$, $\chi_T$ et la relation $P\beta\chi_T = \alpha$.

\begin{enumerate}
\item Exprimer $C_P - C_V$ en fonction de $T$, $V$, $\alpha$ et $\chi_T$ uniquement.
\item La stabilité thermodynamique impose $C_V \geq 0$ et $\chi_T \geq 0$ (vu en cours). En déduire que $\gamma = C_P/C_V \geq 1$. Dans quel cas a-t-on exactement $\gamma = 1$ ?
\item Vérifier que la formule de la question 1 redonne $C_P - C_V = nR$ pour le gaz parfait.
\item Pour l'eau liquide à 4 °C, le coefficient de dilatation $\alpha$ s'annule (anomalie dilatométrique). Que valent alors $C_P - C_V$ et $\gamma$ ? Commenter le cas général des phases condensées.
\end{enumerate}

\begin{indication}
$\left(\partial P/\partial T\right)_V = \beta P$ et $\left(\partial V/\partial T\right)_P = \alpha V$, puis éliminer $\beta$ avec $P\beta\chi_T = \alpha$.
\end{indication}

\begin{solution}
\textbf{Question 1.} $C_P - C_V = T(\beta P)(\alpha V) = TV\alpha\,(\beta P)$, et $\beta P = \alpha/\chi_T$ :
\[
\boxed{C_P - C_V = \frac{TV\alpha^2}{\chi_T}}.
\]

\textbf{Question 2.} Le membre de droite est un carré divisé par $\chi_T \geq 0$ : donc $C_P \geq C_V \geq 0$, d'où $\gamma = C_P/C_V \geq 1$ pour tout système stable — quel que soit le fluide, parfait ou non. L'égalité $\gamma = 1$ exige $\alpha = 0$ : un corps qui ne se dilate pas du tout avec la température.

\textbf{Question 3.} Gaz parfait : $\alpha = 1/T$, $\chi_T = 1/P$, donc $\dfrac{TV\alpha^2}{\chi_T} = \dfrac{TV\,P}{T^2} = \dfrac{PV}{T} = nR$. \checkmark

\textbf{Question 4.} À 4 °C, $\alpha = 0$ donne exactement $C_P = C_V$ et $\gamma = 1$ pour l'eau. Plus généralement, pour les liquides et solides, $\alpha$ est très petit : $C_P - C_V = TV\alpha^2/\chi_T$ est du second ordre en $\alpha$, ce qui justifie l'approximation courante $C_P \approx C_V = C$ pour les phases condensées.
\end{solution}
\end{exo}

\begin{exo}[Équation d'état d'un liquide ou solide]{eos-liquide-solide}
\keywords{équation d'état, liquide, solide, dilatation thermique, compressibilité, phase condensée}

On modélise une phase condensée par $\alpha = \dfrac{1}{V}\left(\dfrac{\partial V}{\partial T}\right)_P$ et $\chi_T = -\dfrac{1}{V}\left(\dfrac{\partial V}{\partial P}\right)_T$ constants. On note $(T_0, P_0, V_0)$ un état de référence.

\begin{enumerate}
\item Intégrer les définitions pour obtenir $V(T,P) = V_0 \exp\!\left[\alpha(T-T_0) - \chi_T(P-P_0)\right]$, puis le développement linéaire à faible écart.
\item On admet $C_P = m c$ constant. Montrer que $S(T,P) = mc\ln(T/T_0) - \alpha V_0 (P-P_0) + \text{cste}$ à l'ordre linéaire en $\alpha$ et $\chi_T$.
\item Comparer au modèle incompressible ($\alpha = \chi_T = 0$) vu en cours : pourquoi ce dernier est-il pathologique ?
\item Application : pour l'eau liquide ($\alpha \approx 2{,}1\times10^{-4}$ K$^{-1}$, $\chi_T \approx 4{,}5\times10^{-10}$ Pa$^{-1}$, $V_0 = 18$ cm$^3$·mol$^{-1}$), estimer la variation de volume molaire entre $P_0 = 1$ bar et $P = 1000$ bar à $T = T_0$.
\end{enumerate}

\begin{indication}
$\mathrm{d}V/V = \alpha\,\mathrm{d}T - \chi_T\,\mathrm{d}P$ à coefficients constants. Pour $S$ : $T\,\mathrm{d}S = C_P\,\mathrm{d}T + h\,\mathrm{d}P$ avec $h = -\alpha VT$.
\end{indication}

\begin{solution}
\textbf{Question 1.} Intégration directe : $\ln(V/V_0) = \alpha(T-T_0) - \chi_T(P-P_0)$. DL : $V \approx V_0[1 + \alpha(T-T_0) - \chi_T(P-P_0)]$.

\textbf{Question 2.} $T\,\mathrm{d}S = mc\,\mathrm{d}T - \alpha V_0\,\mathrm{d}P$ à l'ordre linéaire, d'où $S = mc\ln(T/T_0) - \alpha V_0(P-P_0) + \text{cste}$.

\textbf{Question 3.} Si $\alpha = \chi_T = 0$, alors $V = V_0$ et $S = mc\ln(T/T_0)$ seule : $S = S(U)$, donc $P = T(\partial S/\partial V)_U = 0$ — pression nulle, absurde. Il faut garder $\alpha, \chi_T$ petits mais non nuls.

\textbf{Question 4.} $\Delta V/V_0 \approx -\chi_T \Delta P = -4{,}5\times10^{-10} \times 999\times10^5 \approx -0{,}045$ : environ $-4{,}5\,\%$ de volume, non négligeable à 1000 bar.
\end{solution}
\end{exo}

\begin{exo}[Détente de Joule-Thomson]{joule-thomson}
\keywords{Joule-Thomson, détente isenthalpique, coefficient de Joule-Thomson, van der Waals, liquéfaction}

Une détente de Joule-Thomson est une détente lente (quasi-statique) à enthalpie constante : le fluide passe d'une pression $P_i$ à $P_f < P_i$ à travers un tamis ou une valve.

\begin{enumerate}
\item Montrer que la variation de température s'écrit $\Delta T = \displaystyle\int_{P_i}^{P_f} \mu_{JT}\,\mathrm{d}P$ avec $\mu_{JT} = \left(\dfrac{\partial T}{\partial P}\right)_H$.
\item Établir $\mu_{JT} = \dfrac{1}{C_P}\left[T\left(\dfrac{\partial V}{\partial T}\right)_P - V\right] = \dfrac{V}{C_P}(\alpha T - 1)$.
\item Pour un gaz parfait, que vaut $\mu_{JT}$ ? Retrouver $\Delta T = 0$.
\item Gaz de van der Waals ($C_P$ constant) : montrer que $\mu_{JT} > 0$ pour $T < T_i = 2a/(Rb)$ (température d'inversion). Application : $a = 0{,}14$ Pa·m$^6$·mol$^{-2}$, $b = 3{,}9\times10^{-5}$ m$^3$·mol$^{-1}$ — estimer $T_i$.
\end{enumerate}

\begin{indication}
Enthalpie constante : $\mathrm{d}H = T\,\mathrm{d}S + V\,\mathrm{d}P = 0$. Utiliser $\left(\partial S/\partial P\right)_T = -\left(\partial V/\partial T\right)_P$. Pour VdW : $T_i$ vérifie $\alpha T_i = 1$ avec $\alpha = (V-nb)/(nRTV)$.
\end{indication}

\begin{solution}
\textbf{Question 1.} $\Delta T = T(P_f) - T(P_i) = \int_{P_i}^{P_f} (\partial T/\partial P)_H\,\mathrm{d}P$ par définition de $\mu_{JT}$.

\textbf{Question 2.} Écrire $\mathrm{d}H = T\,\mathrm{d}S + V\,\mathrm{d}P$, et utiliser $\mathrm{d}S = \left(\dfrac{\partial S}{\partial T}\right)_P\mathrm{d}T + \left(\dfrac{\partial S}{\partial P}\right)_T\mathrm{d}P = \dfrac{C_P}{T}\,\mathrm{d}T - \left(\dfrac{\partial V}{\partial T}\right)_P\mathrm{d}P$ (relation de Maxwell). On obtient
\[
\mathrm{d}H = C_P\,\mathrm{d}T + \left[V - T\left(\frac{\partial V}{\partial T}\right)_P\right]\mathrm{d}P.
\]
À $H$ constant, $\mathrm{d}H = 0$ donne directement
\[
\mu_{JT} = \left(\frac{\partial T}{\partial P}\right)_H = \frac{T\left(\partial V/\partial T\right)_P - V}{C_P} = \frac{V(\alpha T - 1)}{C_P}.
\]
Attention à ne pas confondre ce résultat avec le coefficient calorimétrique $h = -TV\alpha$ de l'exercice sur les relations de Clapeyron : $h$ intervient dans $\delta Q = C_P\,\mathrm{d}T + h\,\mathrm{d}P$, alors que $\mu_{JT}$ fait intervenir $H$, et $\mathrm{d}H \neq \delta Q$ dès que $\mathrm{d}P \neq 0$ (on a $\mathrm{d}H = \delta Q + V\,\mathrm{d}P = C_P\mathrm{d}T + (h+V)\mathrm{d}P$).

\textbf{Question 3.} GP : $\alpha = 1/T$, donc $\alpha T - 1 = 0$ : $\mu_{JT} = 0$ et $\Delta T = 0$ (première loi de Joule).

\textbf{Question 4.} Pour VdW, $\alpha T > 1$ sous $T_i = 2a/(Rb) \approx 2\times0{,}14/(8{,}314\times3{,}9\times10^{-5}) \approx 860$ K. En dessous, la détente refroidit ($\mu_{JT} > 0$) : principe de liquéfaction des gaz (Linde, Claude).
\end{solution}
\end{exo}
